\clearpage\pagestyle{otMain}\setlength{\otBlocEspacement}{2.5mm}\fontsize{11}{15}\selectfont\corrigesChapitre{7}{Proportionnalité, rapports et ratios}
\begin{blocLivre}{cor-7-1}\begin{corrige}{1}{Passer à l’unité}\demarche{Calculer le prix d’une unité.}
Une bouteille coûte $9:6=1{,}50$ euro. Dix bouteilles coûtent $10\times1{,}50=\res{15}$ euros. On suppose le prix unitaire constant.
\end{corrige}\end{blocLivre}
\begin{blocLivre}{cor-7-2}\begin{corrige}{2}{Compléter un tableau}\demarche{Préciser le sens et l’unité du coefficient.}
Le coefficient masse vers prix est $\frac72=3{,}5$ euros par kg. Les prix sont donc $5\times3{,}5=\res{17{,}50}$ euros et $8\times3{,}5=\res{28}$ euros.
\end{corrige}\end{blocLivre}
\begin{blocLivre}{cor-7-3}\begin{corrige}{3}{Des produits en croix}\demarche{Écrire l’égalité des rapports avant le produit en croix.}
$\frac6{15}=\frac8x$ donne $6x=15\times8=120$, donc $\res{x=20}$. Vérification : le coefficient vaut $\frac{15}6=2{,}5$ et $8\times2{,}5=20$.
\end{corrige}\end{blocLivre}
\begin{blocLivre}{cor-7-4}\begin{corrige}{4}{Un ratio dans une classe}\demarche{Comparer les deux groupes, pas une partie au total.}
Il y a $28-12=16$ garçons. Le rapport vaut $\frac{12}{16}=\res{\frac34}$. Le ratio filles : garçons est $12:16$, soit $\res{3:4}$ : pour 3 filles, il y a 4 garçons.
\end{corrige}\end{blocLivre}
\begin{blocLivre}{cor-7-5}\begin{corrige}{5}{Lire une échelle}\demarche{Faire d’abord le rapport avec des unités identiques.}
La longueur réelle est $3{,}6\times500=1\,800$ cm, soit $\res{18\ \mathrm m}$. La conversion en mètres vient après le calcul dans la même unité.
\end{corrige}\end{blocLivre}
\begin{blocLivre}{cor-7-6}\begin{corrige}{6}{Une quatrième proportionnelle}\demarche{Comparer passage à l’unité et produits en croix.}
Par unité : $14{,}4:8=1{,}8$ kg, puis $5\times1{,}8=\res9$ kg. Par produits en croix : $\frac8{14{,}4}=\frac5m$ donne $8m=72$, donc $m=9$ kg. Les deux méthodes concordent.
\end{corrige}\end{blocLivre}
\begin{blocLivre}{cor-7-7}\begin{corrige}{7}{Un tarif non proportionnel}\demarche{Tester un doublement.}
Pour 1 h : $4+3=\res7$ euros. Pour 2 h : $4+6=\res{10}$ euros. Doubler le temps ne double pas 7 euros, donc il n’y a $\res{\text{pas proportionnalité}}$. Les frais fixes l’empêchent.
\end{corrige}\end{blocLivre}
\begin{blocLivre}{cor-7-8}\begin{corrige}{8}{Un graphique}\demarche{Vérifier l’alignement avec l’origine puis lire les axes.}
Les points sont sur une droite passant par l’origine : le prix est proportionnel à la masse dans ce modèle. À 3 kg, l’ordonnée est $\res{7{,}50}$ euros. Le coefficient est 2,50 euros par kg.
\label{fig-cor-7-8}\figureGeo[corrige]{proportion-prix}
\end{corrige}\end{blocLivre}
\begin{blocLivre}{cor-7-9}\begin{corrige}{9}{Partager selon un ratio}\demarche{Compter les parts, puis calculer une part.}
On compte $4+5=9$ parts égales ; une part vaut $72:9=8$ euros. Ana reçoit $4\times8=\res{32}$ euros et Ben $5\times8=\res{40}$ euros. Contrôle : $32+40=72$ et $\frac{32}4=\frac{40}5$.
\label{fig-cor-7-9}\figureGeo[corrige]{ratio-72}
\end{corrige}\end{blocLivre}
\begin{blocLivre}{cor-7-10}\begin{corrige}{10}{Trois nombres proportionnels}\demarche{Utiliser le total des parts.}
Le ratio $2:3:7$ compte $2+3+7=12$ parts. Une part vaut $60:12=5$. Donc $\res{a=10}$, $\res{b=15}$ et $\res{c=35}$. On vérifie $\frac{10}2=\frac{15}3=\frac{35}7=5$.
\end{corrige}\end{blocLivre}
\begin{blocLivre}{cor-7-11}\begin{corrige}{11}{Une recette et des lots}\demarche{Calculer la quantité puis arrondir le nombre de lots vers le haut.}
Par personne, il faut $450:6=75$ g. Pour 14 personnes, $75\times14=1\,050$ g. Deux paquets ne donnent que 1 000 g : il faut $\res3$ paquets. La quantité achetée peut dépasser la quantité nécessaire.
\end{corrige}\end{blocLivre}
\begin{blocLivre}{cor-7-12}\begin{corrige}{12}{Une part connue}\demarche{Retrouver la valeur d’une part à partir de la part connue.}
Le deuxième a 3 parts, donc une part vaut $18:3=6$ cartes. Le premier reçoit $4\times6=\res{24}$ cartes, le troisième $2\times6=\res{12}$ cartes. Le total est $9\times6=\res{54}$ cartes.
\end{corrige}\end{blocLivre}
\begin{blocLivre}{cor-7-13}\begin{corrige}{13}{Un arc de cercle}\demarche{Prendre la fraction d’un tour complet.}
La longueur d’un arc est proportionnelle à son angle. La circonférence vaut $2\pi\times6=12\pi$ cm pour $360^\circ$. L’arc mesure donc $12\pi\times\frac{75}{360}=\res{\frac{5\pi}2\ \mathrm{cm}}$, soit environ $\res{7{,}9\ \mathrm{cm}}$.
\label{fig-cor-7-13}\figureGeo[corrige]{arc-75}
\end{corrige}\end{blocLivre}
\begin{blocLivre}{cor-7-14}\begin{corrige}{14}{Un mélange de peinture}\demarche{Distinguer le ratio entre deux couleurs et la fraction du total.}
Le bleu représente 2 parts, donc une part vaut $3:2=1{,}5$ L. Il faut $5\times1{,}5=\res{7{,}5\ \mathrm L}$ de blanc. Le mélange mesure $3+7{,}5=\res{10{,}5\ \mathrm L}$. Le bleu représente $\res{\frac27}$ du mélange, car $2+5=7$ parts.
\end{corrige}\end{blocLivre}
\begin{blocLivre}{cor-7-15}\begin{corrige}{15}{Comparer deux recettes}\demarche{Comparer des quantités rapportées à un litre.}
A : $250:0{,}5=500$ g/L. B : $400:0{,}75\approx533$ g/L. La recette $\res{B}$ est la plus riche en farine. Pour 1,2 L avec B : $\frac{400\times1{,}2}{0{,}75}=\res{640\ \mathrm g}$. Comparer les seules masses de farine n’aurait pas suffi, car les volumes de lait diffèrent.
\end{corrige}\end{blocLivre}
